\chapter{Introduction}
\label{chap:introduction}

\begin{definition}
  Un vecteur de l’espace s’écrit dans la base canonique
  $\vec{v} = a\,\ihat + b\,\jhat + c\,\khat$.
\end{definition}

\begin{theoreme}[Norme euclidienne]
  \label{thm:norme}
  Pour tout $\vec{v} = a\,\ihat + b\,\jhat + c\,\khat$, on a
  \[
    \lVert \vec{v} \rVert = \sqrt{a^2 + b^2 + c^2}.
  \]
\end{theoreme}

\begin{proof}
  Appliquer deux fois le théorème de Pythagore.
\end{proof}

\begin{exemple}
  Si $\vec{v} = 3\,\ihat + 4\,\jhat$, alors $\lVert \vec{v} \rVert = 5$
  (théorème~\ref{thm:norme}).
\end{exemple}

\begin{remarque}
  Ajoutez d’autres chapitres dans le dossier \texttt{chapitres/} et incluez-les avec \verb|\input|.
\end{remarque}
